% Use this file to write code that you do not want in the .sty file, but has to be in the preamble (before \begin{document}).
% Writing in this file in stead of in the preamble will keep the main file more organised and tidy.


%                       Nomenclatures
%_______________________________________________________________
\usepackage[intoc]{nomencl}
\makenomenclature

\renewcommand{\nomname}{%
% Title
%----------------
List of Symbols
%----------------
}
\renewcommand{\nompreamble}{%
% Description
%----------------
The next list describes several symbols that will be later used within the body of the document
%----------------
}
% This code creates the groups
\renewcommand\nomgroup[1]{%
  \item[\bfseries
  \ifstrequal{#1}{A}{Physics constants}{%
  \ifstrequal{#1}{B}{Mathematical constants}{%
  \ifstrequal{#1}{C}{Other symbols}{}}}%
]}
% This will add the units
\newcommand{\nomunit}[1]{%
\renewcommand{\nomentryend}{\\#1}}
%................................................................



%                       Additions for this report
%_______________________________________________________________
\usepackage{multirow}   % Vertically merged table cells.
%................................................................
\usepackage{seqsplit}                   % Break long SMILES strings anywhere.
\setlength{\emergencystretch}{3em}      % Last-resort stretch in narrow columns.
\newcommand{\smiles}[1]{\texttt{\seqsplit{#1}}}
%................................................................

% Tighter float separation: the defaults leave a lot of air around floats
% in a two-column layout, and the report has a hard page budget.
\setlength{\textfloatsep}{10pt plus 2pt minus 2pt}
\setlength{\floatsep}{9pt plus 2pt minus 2pt}
\setlength{\intextsep}{9pt plus 2pt minus 2pt}
\setlength{\dbltextfloatsep}{10pt plus 2pt minus 2pt}
\setlength{\dblfloatsep}{9pt plus 2pt minus 2pt}
